\section{Εξέλιξη Καταλληλότητας και Δυναμική Πληθυσμού}
\label{sec:eval_fitness}

\subsection{Σύγκλιση Καταλληλότητας και Στατιστική Σημαντικότητα}
Η πορεία της εξελικτικής βελτιστοποίησης παρακολουθείται μέσω της συνάρτησης καταλληλότητας $F(\mathbf{x})$, η οποία εκφράζει την καθαρή αεροδυναμική απόδοση (Lift-to-Drag ratio, $L/D$) σε κατάσταση ισόρροπης πτήσης cruise, υπό την προϋπόθεση ικανοποίησης όλων των φυσικών και στατικών περιορισμών. 

Στο Σχήμα~\ref{fig:fitness_evolution} αποτυπώνεται η εξέλιξη της βέλτιστης (Best), μέσης (Mean) και χείριστης (Worst) τιμής καταλληλότητας ανά γενιά κατά μήκος 47 εξελικτικών κύκλων, συνοδευόμενη από ζώνες εμπιστοσύνης $\pm 1\sigma$. Στον Πίνακα~\ref{tab:fitness_evolution} συνοψίζονται οι αντίστοιχες ποσοτικές μετρήσεις ανά κομβικά ορόσημα γενεών.

\begin{figure}[htbp]
    \centering
    \includegraphics[width=0.88\textwidth]{figures/fig1_fitness_evolution.png}
    \caption{Εξέλιξη συνάρτησης καταλληλότητας ($L/D$) κατά μήκος των γενεών: Τροχιές βέλτιστης, μέσης εφικτής και χείριστης λύσης με διάστημα εμπιστοσύνης $\pm 1\sigma$.}
    \label{fig:fitness_evolution}
\end{figure}

../../evaluation/figures/table_fig1_fitness_evolution.tex

Παρατηρείται ότι κατά τις πρώτες 10 γενιές συντελείται απότομη αύξηση της μέσης καταλληλότητας (από $9.37$ σε $11.68$), καθώς ο αλγόριθμος απορρίπτει τα μη εφικτά άτομα και συγκλίνει προς την ευρύτερη περιοχή βέλτιστων αεροδυναμικών διαμορφώσεων. Μετά τη γενιά $t=20$, η βέλτιστη τιμή σταθεροποιείται ασυμπτωτικά στο $L/D \approx 12.28 \pm 0.00$, υποδεικνύοντας επιτυχή προσέγγιση του καθολικού ακροτάτου.

Για την επιβεβαίωση της στατιστικής αναπαραγωγιμότητας του αλγορίθμου και τον αποκλεισμό τυχαίας σύγκλισης, εκτελέστηκαν 5 ανεξάρτητα πειράματα με διαφορετικούς αρχικούς σπόρους ψευδοτυχαίων αριθμών (Seeds 42, 101, 777, 999, 2024). Στο Σχήμα~\ref{fig:multi_seed_ribbon} παρουσιάζεται η συγκεντρωτική κορδέλα διακύμανσης (statistical ribbon) $\pm 1\sigma$, ενώ στον Πίνακα~\ref{tab:multi_seed_convergence} παρατίθενται οι επιμέρους τιμές.

\begin{figure}[htbp]
    \centering
    \includegraphics[width=0.88\textwidth]{figures/fig5_fitness_multi_seed_ribbon.png}
    \caption{Στατιστική επαλήθευση σύγκλισης πολλαπλών ανεξάρτητων εκτελέσεων (5 distinct seeds): Συγκεντρωτική κορδέλα $\pm 1\sigma$ που αποδεικνύει την εύρωστη σύγκλιση στην ίδια βέλτιστη τιμή $L/D \approx 12.27$ ανεξαρτήτως αρχικοποίησης.}
    \label{fig:multi_seed_ribbon}
\end{figure}

../../evaluation/figures/table_fig5_multi_seed_ribbon.tex

Όπως καταδεικνύει ο Πίνακας~\ref{tab:multi_seed_convergence}, παρά τη σημαντική απόκλιση της αρχικής καταλληλότητας στη γενιά $t=1$ ($7.70$ έως $10.63$, μέση τιμή $9.31 \pm 1.03$), όλες οι ανεξάρτητες εκτελέσεις συγκλίνουν στη γενιά $t=45$ σε τελική βέλτιστη τιμή $12.27 \pm 0.02$, επιβεβαιώνοντας την υψηλή αξιοπιστία της διαδικασίας.

\subsection{Σύγκριση Μονής Νήσου έναντι Κατανεμημένης Πολυ-Νησιωτικής Τοπολογίας}
Για την αξιολόγηση της χρησιμότητας της κατανεμημένης αρχιτεκτονικής νησιών (Multi-Island Ring Topology), συγκρίθηκε η επίδοση ενός μονολιθικού πληθυσμού $P=200$ (Single Island) έναντι δύο διασυνδεδεμένων νησιών μεγέθους $P=100$ έκαστο, με κυκλική ανταλλαγή κορυφαίων μεταναστών κάθε $\tau = 5$ γενιές.

Στο Σχήμα~\ref{fig:single_vs_multi_island} και στον Πίνακα~\ref{tab:single_vs_multi_island} παρουσιάζεται η συγκριτική δυναμική σύγκλισης.

\begin{figure}[htbp]
    \centering
    \includegraphics[width=0.88\textwidth]{figures/fig1_single_vs_multi_island.png}
    \caption{Σύγκριση ταχύτητας και ποιότητας σύγκλισης: Μονή Νήσος (Single Island) έναντι Κατανεμημένης Πολυ-Νησιωτικής Τοπολογίας Δακτυλίου (Multi-Island Ring).}
    \label{fig:single_vs_multi_island}
\end{figure}

../../evaluation/figures/table_fig1_single_vs_multi_island.tex

Η πολυ-νησιωτική τοπολογία επιτυγχάνει ταχύτερη ανακάλυψη ανώτερων διαμορφώσεων στις πρώιμες γενιές ($\Delta = +0.14$ στο $t=1$, $+0.10$ στο $t=5$), διατηρώντας ανεξάρτητους δρόμους αναζήτησης σε κάθε υπο-πληθυσμό και αποφεύγοντας την πρόωρη παγίδευση σε τοπικά βέλτιστα.

\subsection{Δυναμική Ποικιλομορφίας Πληθυσμού και Εντροπία Shannon}
Η διατήρηση της γενετικής ποικιλομορφίας αποτελεί κρίσιμο παράγοντα για την αποφυγή πρόωρης γενετικής παρέκκλισης (genetic drift). Για την ποσοτικοποίησή της υπολογίστηκαν δύο μετρικές κατά μήκος των γενεών:
\begin{enumerate}
    \item \textbf{Κανονικοποιημένη Εντροπία Shannon $H(t)$:} Υπολογιζόμενη επί των ιστογραμμάτων κατανομής των 10 γονιδίων:
    \begin{equation}
        H(t) = -\frac{1}{10 \ln(B)} \sum_{j=1}^{10} \sum_{b=1}^{B} p_{j,b} \ln(p_{j,b})
    \end{equation}
    όπου $B=20$ είναι ο αριθμός των διακριτών κάδων (bins) και $p_{j,b}$ η σχετική συχνότητα του γονιδίου $j$ στον κάδο $b$.
    \item \textbf{Μέση Γονιδιακή Διασπορά (Mean Gene Variance $\sigma_g^2$):} Η μέση διακύμανση των κανονικοποιημένων γονιδίων $\mathbf{x} \in [0, 1]^{10}$.
\end{enumerate}

Στο Σχήμα~\ref{fig:population_diversity} και στον Πίνακα~\ref{tab:population_diversity} απεικονίζεται η εξέλιξη της ποικιλομορφίας.

\begin{figure}[htbp]
    \centering
    \includegraphics[width=0.95\textwidth]{figures/fig1_population_diversity.png}
    \caption{Δυναμική ποικιλομορφίας πληθυσμού: (α) Κανονικοποιημένη Εντροπία Shannon $H(t)$ και (β) Μέση Γονιδιακή Διασπορά για Μονή Νήσο και Κατανεμημένο Δακτύλιο Νήσων.}
    \label{fig:population_diversity}
\end{figure}

../../evaluation/figures/table_fig1_population_diversity.tex

Η εντροπία ξεκινά κοντά στη μέγιστη τιμή ($0.98$ στο $t=1$) λόγω της ομοιόμορφης τυχαίας αρχικοποίησης, μειώνεται ομαλά κατά τη φάση εκμετάλλευσης (exploitation) και σταθεροποιείται στο $H \approx 0.29$ κατά τις τελευταίες γενιές. Η διατήρηση διασποράς $\sigma_g^2 \approx 0.006$ αποδεικνύει ότι ο πληθυσμός διατηρεί επαρκή γενετική ευελιξία χωρίς να καταρρέει σε ένα και μόνο μονοδιάστατο σημείο.

\subsection{Προβολή Αναζήτησης σε Πολλαπλότητα PCA}
Για να γίνει κατανοητή η γεωμετρία της σύγκλισης στον 10-διάστατο χώρο γονιδιώματος, εφαρμόστηκε Ανάλυση Κύριων Συνιστωσών (Principal Component Analysis - PCA) στο σύνολο των αξιολογηθέντων ατόμων.

Στο Σχήμα~\ref{fig:pca_manifold} παρουσιάζεται η προβολή των πληθυσμών των γενεών στο 2D επίπεδο των δύο πρώτων κύριων συνιστωσών (PC1 και PC2). Στον Πίνακα~\ref{tab:pca_loadings} αναλύονται τα βάρη φόρτισης (loadings) των παραμέτρων σχεδιασμού στις δύο πρώτες συνιστώσες, οι οποίες εξηγούν συνδυαστικά το $60.4\%$ της συνολικής διακύμανσης του σχεδιαστικού χώρου.

\begin{figure}[htbp]
    \centering
    \includegraphics[width=0.90\textwidth]{figures/fig1_pca_evolution_manifold.png}
    \caption{Προβολή εξελικτικής πορείας στην πολλαπλότητα 2D PCA: Χαρτογράφηση του 10D γενετικού χώρου όπου καταγράφεται η μετάβαση του πληθυσμού από την ευρεία εξερεύνηση στην πυκνή συσσώρευση επί της αεροδυναμικά βέλτιστης περιοχής.}
    \label{fig:pca_manifold}
\end{figure}

../../evaluation/figures/table_fig1_pca_loadings.tex

Όπως καταγράφεται στον Πίνακα~\ref{tab:pca_loadings}, η κύρια συνιστώσα PC1 (εξηγεί το $42.5\%$) καθοδηγείται κατά κύριο λόγο από τη δίεδρο ($\Gamma$, $+0.501$), τη γωνία συστροφής ($\theta_{\text{twist}}$, $+0.464$) και τη χορδή ρίζας ($c_{\text{root}}$, $+0.395$), οι οποίες καθορίζουν τη διαμήκη ευστάθεια και την αποφυγή αποκόλλησης. Η PC2 (εξηγεί το $17.9\%$) κυριαρχείται από τη χορδή ακροπτερυγίου ($c_{\text{tip}}$, $+0.762$) που ελέγχει το λόγο επιμήκυνσης (Aspect Ratio) και την οπισθέλκουσα επαγωγής.
